% Candidate compact structural upgrade for the main Erdős 593 manuscript.
% Intended placement immediately after the canonical atom normal form.

\begin{proposition}[Universal canonical one-point factorization]
\label{proposition-universal-canonical-one-point-factorization}
Let $F$ be a finite reduced obligatory triple system.  For hyperedges $e,f$
in the same connected component, write $e\sim f$ when for every point
$p\in V(F)$ the hyperedge-nodes $e$ and $f$ remain in the same connected
component of $I(F)-p$.

Then $\sim$ is an equivalence relation, and its equivalence classes are
exactly the hyperedge sets of the canonical atoms.  Moreover a point belongs
to at least two canonical atoms if and only if deleting that point increases
the number of Levi components containing hyperedge-nodes.

Consequently the canonical atom decomposition is the unique finest
decomposition of $E(F)$ into connected supported pieces that meet pairwise in
at most one point and whose piece--shared-point incidence graph is a forest.
Every other such decomposition is obtained by grouping whole canonical atoms.
\end{proposition}

\begin{proof}
Each canonical atom is one-point indecomposable: it is either one triple or
$J^+$ with $J$ finite, simple, bipartite and $2$-connected.  Hence deleting one
point never separates two hyperedge-nodes belonging to the same atom.

If two hyperedges belong to distinct atoms in one connected component, the
unique path between their atom nodes in the canonical atom--shared-point
forest contains a shared point $p$.  Removing $p$ separates the two branches,
and exact reconstruction shows that the corresponding hyperedge-nodes lie in
different components of $I(F)-p$.  This proves the first assertion and also
shows that every shared point is detected by deletion.  Conversely, if a
point lies in only one atom, one-point indecomposability keeps that atom's
hyperedge support connected after deletion, while all attachments through the
other shared points remain unchanged.  Thus such a point is not a separator.

Finally consider any supported one-point forest decomposition.  If it split a
canonical atom between two pieces, the unique path between those pieces in the
piece--shared-point forest would contain a shared point whose deletion
separates two hyperedges of that atom, contradicting the first part.  Hence
every canonical atom is contained in one piece, proving the universal
refinement statement.
\end{proof}

\begin{corollary}[Decomposition lattice]
\label{corollary-canonical-decomposition-lattice}
Let $B(F)$ be the graph whose vertices are the canonical atoms, with two atoms
adjacent when they share a point.  Then $B(F)$ is a block graph; its maximal
cliques of size at least two are precisely the sets of atoms through a fixed
shared point.

The supported one-point forest decompositions of $F$, ordered by refinement,
are in canonical bijection with the set partitions of $V(B(F))$ in which
each block induces a connected subgraph.  Equivalently, the decomposition
poset is the bond lattice of $B(F)$.
\end{corollary}

\begin{proof}
A shared point gives a clique of incident atoms.  A cycle in $B(F)$ using
more than one shared point would lift to a cycle in the bipartite
atom--shared-point incidence forest, so every graph block of $B(F)$ is one of
these cliques.

By the proposition, every supported one-point forest decomposition groups
whole canonical atoms.  A piece is connected exactly when its atom set is
connected in $B(F)$.  Conversely, unions over the blocks of any connected set
partition are connected; two such unions cannot meet in two points, and a
cycle in their piece--shared-point incidence graph would lift to a cycle in
the canonical atom forest.  Thus they form a supported one-point forest
decomposition.  The constructions are inverse and preserve refinement.
\end{proof}
